\subsection{The case of normed spaces}

Let \(\mathbf{E}\) be a normed space. The strong topology on the dual \(\mathbf{E}'\) of \(\mathbf{E}\) is defined by the norm

\[
\left\| x ^ {\prime} \right\| = \sup _ {x \in \mathrm{E}, \| x \| \leqslant 1} \left| \langle x, x ^ {\prime} \rangle \right|,\tag{1}
\]

and the strong dual of E is a Banach space (III, p.~24, cor. 2). Then the bidual E'' of E is also a Banach space, for the norm defined by

\[
\left\| x ^ {\prime \prime} \right\| = \sup _ {x ^ {\prime} \in \mathrm{E} ^ {\prime}, \left\| x ^ {\prime} \right\| \leqslant 1} \left| \left\langle x ^ {\prime}, x ^ {\prime \prime} \right\rangle \right|.\tag{2}
\]

By prop. 8, (i) of IV, p.~7, the canonical linear mapping \(c_{\mathrm{E}}: \mathrm{E} \to \mathrm{E}''\) is an isometry. Henceforth, we shall identify E with a normed subspace of its bidual \(\mathrm{E}''\) .

PROPOSITION 5. --- Let E be a normed space, E' its dual and E'' its bidual. The (closed) unit ball in E'' is the closure of the unit ball B in E for the weak topology \(\sigma(\mathrm{E}'', E')\) . By formulas (1) and (2), the unit ball in E'' is the bipolar B°° of B. Prop. 5 then follows from the theorem of bipolars (II, p.~45, cor. 3).

Remark. --- A Banach space E is closed in its bidual E'' for the strong topology, but is dense for the weak topology (prop. 5).

In order that a normed space be reflexive, it is necessary and sufficient that it is semi-reflexive; for, the initial topology of E is always induced by the strong topology of \(E''\) . Th. 1 (IV, p.~15) then implies the following result:

PROPOSITION 6. --- In order that a normed space E be reflexive, it is necessary and sufficient that the unit ball in E be compact for the weakened topology \(\sigma(E, E')\) .

We observe that a reflexive normed space is complete hence a Banach space, and that its dual is a reflexive Banach space by prop. 4 of IV, p.~16.

PROPOSITION 7. --- Let E be a reflexive Banach space and M a closed vector subspace of E. Then M and E/M are reflexive Banach spaces.

Let \(E'\) be the dual of E and \(M^{\circ}\) the orthogonal of M in \(E'\) . As a normed space, we can identify the space \(E'/M^{\circ}\) with the dual \(M'\) of M (IV, p.~9, prop. 10). Since M is semi-reflexive (IV, p.~16, corollary), it is reflexive, hence so is \(E'/M^{\circ}\) ; similarly \(M^{\circ}\) is reflexive, as also its dual \(E/M^{\circ\circ} = E/M\) .

Examples. --- 1) Let \(\ell^{\infty}(\mathbf{N})\) denote the Banach space of bounded sequences \(\mathbf{x} = (x_n)_{n\in \mathbf{N}}\)

of scalars, with the norm

\[
\| \mathbf {x} \| = \sup _ {n \in \mathbf {N}} | x _ {n} | \quad (\mathrm{I}, \text {p.} 4).\tag{3}
\]

Let \(c_{0}(\mathbf{N})\) be the closed vector subspace of \(\ell^{\infty}(\mathbf{N})\) consisting of sequences tending to 0. Finally, let \(\ell^{1}(\mathbf{N})\) be the vector space of summable sequences, endowed with the norm

\[
\left\| \mathbf {x} \right\| _ {1} = \sum_ {n \in \mathbf {N}} \left| x _ {n} \right|.\tag{4}
\]

We can show (IV, p.~47, exerc. 1) that the dual of \(c_0(\mathbf{N})\) can be identified with \(\ell^1 (\mathbf{N})\) in such a way that we have

\[
\langle \mathbf {x}, \mathbf {x} ^ {\prime} \rangle = \sum_ {n \in \mathbf {N}} x _ {n} x _ {n} ^ {\prime}\tag{5}
\]

for all \(\mathbf{x} \in c_0(\mathbf{N})\) and \(\mathbf{x}' \in \ell^1(\mathbf{N})\) . Similarly the dual of \(\ell^1(\mathbf{N})\) can be identified with \(\ell^\infty(\mathbf{N})\) in such a way that we have the relation (5) for all \(\mathbf{x} \in \ell^1(\mathbf{N})\) and all \(\mathbf{x}' \in \ell^\infty(\mathbf{N})\) . Hence \(\ell^\infty(\mathbf{N})\) is the bidual of \(c_0(\mathbf{N})\) , and this latter space is not reflexive.

* 2) Every hilbertian space is a reflexive Banach space (V, p.~17). *

* 3) Let X be a Hausdorff topological space and \(\mu\) a complex measure on X. For every real number \(p > 1\) , the Banach space \(\mathrm{L}^p (\mathrm{X},\mu)\) is reflexive, and its dual can be identified with \(\mathrm{L}^q (\mathrm{X},\mu)\) with \(p^{-1} + q^{-1} = 1\) (INT, V, 2nd edition, § 5, No.~8 and IX, § 1, No.~10). *

